%% ===========================================================================
%%  MULTIVARIATE DATA ANALYSIS — A Computational Companion with R
%%  MS_LJMR edition. Revised course code (Soberon / Jimenez / Cobos) integrated
%%  with the Machine-Learning course of Luis J. Madrigal-Roca.
%%
%%  Single-file LaTeX book. Compile:  pdflatex x3  (from the book/ directory)
%%  Figures live in ../figs/ (see \graphicspath).
%% ===========================================================================
\documentclass[11pt,oneside]{book}

\usepackage[T1]{fontenc}
\usepackage[utf8]{inputenc}
\usepackage{lmodern}
\usepackage[scaled=0.92]{helvet}      % sans for headings
\usepackage{microtype}
\usepackage[a4paper,margin=1in,headheight=15pt]{geometry}
\usepackage{amsmath,amssymb,bm}
\usepackage{graphicx}
\graphicspath{{../figs/}}
\usepackage{booktabs,array}
\usepackage{caption}
\usepackage{xcolor}
\usepackage{titlesec}
\usepackage{fancyhdr}
\usepackage{enumitem}
\usepackage{listings}
\usepackage[most,breakable]{tcolorbox}
\usepackage[hidelinks]{hyperref}

%% ---- Wiley-inspired palette ----------------------------------------------
\definecolor{wnavy}{HTML}{1F3B73}
\definecolor{wwine}{HTML}{8C2D3A}
\definecolor{wgold}{HTML}{C8A23A}
\definecolor{wteal}{HTML}{2A7F7F}
\definecolor{wgray}{HTML}{5B5B5B}
\definecolor{codebg}{HTML}{F5F6F8}
\definecolor{outbg}{HTML}{ECECEC}

\hypersetup{colorlinks=true,linkcolor=wnavy,citecolor=wwine,urlcolor=wteal}
\captionsetup{font=small,labelfont={bf,color=wnavy}}

%% ---- Chapter / section styling (Wiley-like colored numerals) --------------
\titleformat{\chapter}[display]
  {\sffamily\bfseries}
  {\raggedleft\color{wgray}\fontsize{60}{60}\selectfont\thechapter}
  {6pt}
  {\raggedright\color{wnavy}\Huge}
  [\vspace{2pt}{\color{wgold}\titlerule[2pt]}]
\titlespacing*{\chapter}{0pt}{10pt}{28pt}

\titleformat{\section}
  {\sffamily\Large\bfseries\color{wnavy}}{\thesection}{0.6em}{}
\titleformat{\subsection}
  {\sffamily\large\bfseries\color{wwine}}{\thesubsection}{0.6em}{}

%% ---- Running headers -------------------------------------------------------
\pagestyle{fancy}\fancyhf{}
\renewcommand{\headrulewidth}{0.4pt}
\renewcommand{\chaptermark}[1]{\markboth{\thechapter.\ #1}{}}
\fancyhead[L]{\sffamily\small\color{wgray}\leftmark}
\fancyhead[R]{\sffamily\small\color{wgray}\thepage}

%% ---- R code listing style --------------------------------------------------
\lstdefinelanguage{Rlang}{
  keywords={if,else,for,while,function,return,in,TRUE,FALSE,NULL,NA,library,
    require,source,set.seed,data,matrix,c,list,scale,prcomp,princomp,eigen,svd,
    kmeans,hclust,dist,cmdscale,lda,glm,cor,cov,table,factor,predict,rpart,
    randomForest,train,nnet,cutree},
  sensitive=true,
  morecomment=[l]{\#},
  morestring=[b]",
  morestring=[b]'
}
\lstset{
  language=Rlang,
  basicstyle=\ttfamily\footnotesize,
  backgroundcolor=\color{codebg},
  keywordstyle=\color{wnavy}\bfseries,
  commentstyle=\color{wteal}\itshape,
  stringstyle=\color{wwine},
  numbers=left,numberstyle=\tiny\color{wgray},numbersep=8pt,
  showstringspaces=false,
  breaklines=true,
  frame=leftline,framerule=2pt,rulecolor=\color{wgold},
  xleftmargin=14pt,aboveskip=10pt,belowskip=6pt,
  columns=fullflexible,keepspaces=true
}
%% console-output style
\lstdefinestyle{rout}{
  language={},basicstyle=\ttfamily\footnotesize\color{black!85},
  backgroundcolor=\color{outbg},numbers=none,frame=none,
  xleftmargin=14pt,breaklines=true,aboveskip=4pt,belowskip=10pt,
  columns=fullflexible,keepspaces=true
}

%% ---- Callout boxes ---------------------------------------------------------
\newtcolorbox{keybox}[1][]{breakable,enhanced,colback=wnavy!6,colframe=wnavy,
  boxrule=0.8pt,left=8pt,right=8pt,top=6pt,bottom=6pt,arc=2pt,
  fonttitle=\sffamily\bfseries\color{white},coltitle=white,
  colbacktitle=wnavy,title={Key idea},#1}
\newtcolorbox{pitfall}[1][]{breakable,enhanced,colback=wwine!6,colframe=wwine,
  boxrule=0.8pt,left=8pt,right=8pt,top=6pt,bottom=6pt,arc=2pt,
  fonttitle=\sffamily\bfseries\color{white},coltitle=white,
  colbacktitle=wwine,title={Pitfall},#1}
\newtcolorbox{notebox}[1][]{breakable,enhanced,colback=wteal!6,colframe=wteal,
  boxrule=0.8pt,left=8pt,right=8pt,top=6pt,bottom=6pt,arc=2pt,
  fonttitle=\sffamily\bfseries\color{white},coltitle=white,
  colbacktitle=wteal,title={From the ML course},#1}

\newcommand{\code}[1]{\texttt{#1}}
\newcommand{\R}{\textsf{R}}
\newcommand{\pkg}[1]{\textsf{#1}}

\setlength{\parskip}{0.4em}
\setlength{\parindent}{0pt}

% ===========================================================================
\begin{document}
% ===========================================================================

%% ---------------------------- TITLE PAGE ----------------------------------
\begin{titlepage}
\thispagestyle{empty}
\newgeometry{margin=0pt}
\begin{tikzpicture}[remember picture,overlay]
  \fill[wnavy] (current page.north west) rectangle ([yshift=-6.2cm]current page.north east);
  \fill[wgold] ([yshift=-6.2cm]current page.north west) rectangle ([yshift=-6.55cm]current page.north east);
\end{tikzpicture}
\vspace*{1.2cm}
\begin{flushleft}\hspace{1.6cm}
\begin{minipage}{0.86\textwidth}
  \color{white}\sffamily
  {\fontsize{30}{34}\selectfont\bfseries Multivariate Data Analysis}\\[6pt]
  {\LARGE A Computational Companion with \R}\\[10pt]
  {\large\itshape from ordination to machine learning}
\end{minipage}
\end{flushleft}
\vspace{6.4cm}
\vspace{6.4cm}
\begin{flushright}
	\begin{minipage}{0.8\textwidth}\raggedleft\sffamily
		{\Large\bfseries\color{wnavy} Jorge Sober\'on \;$\cdot$\; Laura Jim\'enez \;$\cdot$\; Marlon E. Cobos}\\[4pt]
		{\large\color{wgray} revised and extended by}\\[2pt]
		{\Large\bfseries\color{wwine} Luis J. Madrigal-Roca}\\[16pt]
		{\color{wgray}\rule{0.5\textwidth}{0.8pt}}\\[6pt]
		{\large\color{wgray}\textsc{MS\_LJMR edition}\;\textbullet\; Fall 2026}\\
		{\small\color{wgray} The University of Kansas}
	\end{minipage}\hspace*{2cm} % <-- Adjust this value to increase or decrease the right margin
\end{flushright}
\vfill
\begin{center}\sffamily\small\color{wgray}
  Machine-learning sessions integrated from the author's graduate ML course.
\end{center}
\vspace{1cm}
\restoregeometry
\end{titlepage}

%% ---------------------------- COLOPHON ------------------------------------
\thispagestyle{empty}
\vspace*{\fill}
\noindent{\sffamily\small\color{wgray}
\textbf{Multivariate Data Analysis: A Computational Companion with \R} — MS\_LJMR edition.\\[4pt]
This companion accompanies BIOL 943 (Multivariate Methods) at the University of Kansas.
The original teaching scripts were written by Jorge Sober\'on, Laura Jim\'enez and
Marlon E. Cobos. This edition revises them into a portable, reproducible code base
(\code{MS\_LJMR/}) and integrates four machine-learning sessions adapted from the
graduate ML course of Luis J. Madrigal-Roca.\\[6pt]
\emph{Reproducibility note.} Every figure and numerical result in this book was produced
from the accompanying code with fixed random seeds. Where the original course data sets
were unavailable, a seeded, structure-faithful simulation was used so that all listings
run unmodified; dropping the real files into \code{MS\_LJMR/data/} restores the classroom
examples automatically.\\[6pt]
Text and code released under GPL-2, matching the original course material.}
\vspace*{6cm}

%% ---------------------------- CONTENTS ------------------------------------
\tableofcontents

%% ---------------------------- PREFACE -------------------------------------
\chapter*{Preface}
\addcontentsline{toc}{chapter}{Preface}
\markboth{Preface}{}

This book has one governing idea: \emph{almost everything in multivariate analysis
is a statement about a data matrix and the geometry of its rows and columns}. A table
of $n$ objects measured on $p$ variables can be read two ways — the \emph{Q view}
(relationships among objects) and the \emph{R view} (relationships among variables) —
and the great methods of the field are, at bottom, different ways of decomposing that
one matrix. Principal components, correspondence analysis, discriminant analysis and
even modern machine-learning embeddings all fall out of the same small set of linear-
algebra moves: centering, scaling, distance, and eigen/singular-value decomposition.

The material grew out of Dr. Jorge Sober\'on's Multivariate Methods course. The revised
Fall-2026 schedule folds Principal Coordinate Analysis into the Principal Components
sequence, retires the standalone visualization sessions, and adds four sessions on
machine learning. This edition follows that schedule. It also rebuilds the course's \R{}
scripts into a single, portable package, \code{MS\_LJMR/}, that runs on Windows, macOS,
Linux and in batch mode without editing a single path.

Each chapter follows the same rhythm: the \emph{idea} and its mathematics; the \emph{code}
that expresses it; the \emph{output} it produces; and an \emph{interpretation} that says
what the numbers and pictures mean. Boxes flag the one thing worth remembering
(\textcolor{wnavy}{\textbf{Key idea}}), the trap that catches everyone
(\textcolor{wwine}{\textbf{Pitfall}}), and the places where the machine-learning course
connects back to the classical theory (\textcolor{wteal}{\textbf{From the ML course}}).

%\begin{flushright}\itshape L.\,J.\,M.-R., Lawrence, Kansas\end{flushright}

\mainmatter

% ###########################################################################
\part{Foundations}
% ###########################################################################

% ===========================================================================
\chapter{The Multivariate Data Matrix}
% ===========================================================================
\section{Two views of one table}
A multivariate data set is a matrix $\mathbf{X}$ with $n$ rows (objects, sites,
specimens, samples) and $p$ columns (variables, traits, features). Nearly every
technique in this book answers one of two questions about $\mathbf{X}$, historically distinguished in ecological literature as the Q and R modes \cite{Legendre_2012}:

\begin{itemize}[leftmargin=1.4em]
  \item \textbf{Q~mode.} How similar are the \emph{objects} (rows)? This mode focuses on comparing the sample units, calling for a \emph{dissimilarity} or distance between rows. It is the natural view when seeking to cluster or ordinate sites based on their species composition or objects based on their properties.
  \item \textbf{R~mode.} How associated are the \emph{variables} (columns)? This mode focuses on the relationships among descriptors, calling for a \emph{covariance} or correlation between columns. It is the natural view when aiming to reduce the number of variables by finding latent structures or underlying gradients.
\end{itemize}

\begin{keybox}
The non-zero eigenvalues of $\mathbf{X}\mathbf{X}^{\!\top}$ (an $n\times n$, Q-mode matrix)
and of $\mathbf{X}^{\!\top}\mathbf{X}$ (a $p\times p$, R-mode matrix) are \emph{identical}.
This single fact is why the same decomposition can be read as an ordination of objects or
of variables, and it underlies PCA, correspondence analysis and biplots alike.
\end{keybox}

\section{Centering, scaling, and why units matter}
Raw variables usually live on incompatible scales (millimetres vs.\ millimetres of
rainfall). Two operations put them on common footing. \emph{Centering} subtracts each
column mean, moving the origin to the centroid; \emph{scaling} divides by each column
standard deviation, making variances comparable. The combined $z$-transform,
\[
  z_{ij}=\frac{x_{ij}-\bar x_j}{s_j},
\]
is the default preprocessing for correlation-based methods. In \R{} it is one call,
\code{scale()}, and the covariance of the scaled matrix \emph{is} the correlation matrix
of the original.

\section{Distance and dissimilarity}
For the Q view we need a distance $d(i,h)$ between rows $i$ and $h$. The workhorse is
Euclidean distance on standardized data,
\[
  d(i,h)=\sqrt{\textstyle\sum_{j=1}^{p}(z_{ij}-z_{hj})^2},
\]
but Euclidean distance often fails in ecological contexts, particularly due to the ``double-zero problem'' where joint absences do not imply similarity \cite{Legendre_2012}. Therefore, count and presence/absence data demand other metrics such as the Bray--Curtis dissimilarity \cite{Bray_1957}, Jaccard, or $\chi^2$ distance. When working with mixed numeric and categorical data, Gower's general coefficient of similarity \cite{Gower_1971} is required because it gracefully handles missing values and variables on completely different scales by standardizing each pairwise comparison by the variable's range. Ultimately, choosing the distance is the single most consequential modelling decision in ordination — it decides what ``close'' means before any algorithm runs.

\begin{lstlisting}
## Q-mode dissimilarity and R-mode correlation on Anderson's iris
iris <- get_data("iris")
Xs   <- scale(iris[, 1:4])          # z-transform (mean 0, sd 1)
D    <- dist(Xs, method = "euclidean")   # 150 x 150 distances (Q view)
R    <- cor(iris[, 1:4])            # 4 x 4 correlations       (R view)
round(R, 2)
\end{lstlisting}

\begin{lstlisting}[style=rout]
             Sepal.Length Sepal.Width Petal.Length Petal.Width
Sepal.Length         1.00       -0.12         0.87        0.82
Sepal.Width         -0.12        1.00        -0.43       -0.37
Petal.Length         0.87       -0.43         1.00        0.96
Petal.Width          0.82       -0.37         0.96        1.00
\end{lstlisting}

The correlation block (Figure~\ref{fig:corr}) shows the two petal variables almost
perfectly correlated ($r=0.96$) and sepal width standing apart. In the Q view
(Figure~\ref{fig:diss}), ordering the 150 objects by species reveals a clear block
structure: \emph{setosa} forms a tight, well-separated block, while \emph{versicolor}
and \emph{virginica} bleed into each other — a foreshadowing of every clustering and
classification result later in the book.

\begin{figure}[htbp]\centering
  \includegraphics[width=0.62\linewidth]{corr_heatmap.png}
  \caption{R-mode correlation matrix of the four iris measurements. Petal length and
  width are nearly collinear; sepal width is weakly and negatively related to the rest.}
  \label{fig:corr}
\end{figure}

\begin{figure}[htbp]\centering
  \includegraphics[width=0.7\linewidth]{diss_heatmap.png}
  \caption{Q-mode Euclidean dissimilarity among the 150 iris objects, ordered by
  species. Dark blocks on the diagonal are within-species similarity; \emph{setosa}
  (first block) is sharply separated from the other two.}
  \label{fig:diss}
\end{figure}

\begin{pitfall}
Distances computed on \emph{unstandardized} variables are dominated by whichever
variable happens to have the largest numerical range. Always decide, deliberately,
whether to scale — and remember that Euclidean distance is rarely appropriate for raw
species counts.
\end{pitfall}

% ===========================================================================
\chapter{The Linear Algebra of Ordination}
% ===========================================================================
\section{Transformations as matrices}
Every linear operation on data — rotation, stretching, projection — is a matrix
multiplication $\mathbf{Y}=\mathbf{U}\mathbf{X}$. A symmetric positive-definite matrix
such as a covariance matrix has a \emph{spectral decomposition}
$\mathbf{\Sigma}=\mathbf{V}\mathbf{\Lambda}\mathbf{V}^{\!\top}$, where the columns of
$\mathbf{V}$ (the eigenvectors) are orthogonal directions and the diagonal of
$\mathbf{\Lambda}$ (the eigenvalues) gives the variance along each. Geometrically, the
eigenvectors are the axes of the data's dispersion ellipsoid and the eigenvalues are
the squared semi-axis lengths. This is the engine of PCA.

\section{The singular value decomposition}
The most general and numerically stable move in linear algebra is the \emph{singular value decomposition}
(SVD). Any $n\times p$ matrix factors as
\[
  \mathbf{X}=\mathbf{U}\,\mathbf{D}\,\mathbf{V}^{\!\top},
\]
with $\mathbf{U}$ ($n\times r$) and $\mathbf{V}$ ($p\times r$) containing the left and right singular vectors (both orthonormal), and
$\mathbf{D}$ being the diagonal matrix of singular values $d_1\ge d_2\ge\cdots\ge 0$. 

The remarkable power of SVD lies in its truncation properties. By retaining only the first $k$ terms, we construct
$\mathbf{X}_k=\sum_{i=1}^{k} d_i\,\mathbf{u}_i\mathbf{v}_i^{\!\top}$, which the Eckart--Young--Mirsky theorem guarantees is the absolute best
rank-$k$ approximation of the original matrix $\mathbf{X}$ in the least-squares sense (i.e., minimizing the Frobenius norm of the difference) \cite{Eckart_1936}. Consequently, SVD is simultaneously
the basis of PCA, correspondence analysis, latent semantic analysis, and the low-rank
compression we illustrate next.

\begin{lstlisting}
## Low-rank reconstruction of a grayscale wing image (Limenitis)
G <- get_data("limenitis")          # 69 x 100 intensity matrix
s <- svd(G)
approx <- function(k) s$u[,1:k] %*% diag(s$d[1:k]) %*% t(s$v[,1:k])
err <- sapply(c(1,5,10,50),
              function(k) norm(G - approx(k), "F") / norm(G, "F"))
round(err, 3)
\end{lstlisting}

\begin{lstlisting}[style=rout]
[1] 0.197 0.087 0.077 0.020     # relative Frobenius error at rank 1, 5, 10, 50
\end{lstlisting}

\begin{figure}[htbp]\centering
  \includegraphics[width=\linewidth]{svd_compression.png}
  \caption{SVD compression of a wing image. A rank-1 layer captures the gross shading;
  by rank 5 the structure is essentially recovered, and the singular-value spectrum
  (below) explains why.}
  \label{fig:svdc}
\end{figure}

\begin{figure}[htbp]\centering
  \includegraphics[width=0.56\linewidth]{svd_scree.png}
  \caption{Singular-value spectrum (log scale). The rapid decay means a handful of
  components carry most of the matrix's information — the premise of every dimension-
  reduction method in Part II.}
  \label{fig:svds}
\end{figure}

\begin{keybox}
PCA is SVD applied to a centered (and usually scaled) data matrix. Correspondence
analysis is SVD applied to a $\chi^2$-standardized contingency table. Latent semantic
analysis is SVD applied to a term--document matrix. Learn SVD once; recognise it
everywhere.
\end{keybox}

% ###########################################################################
\part{Ordination and Dimension Reduction}
% ###########################################################################

% ===========================================================================
\chapter{Principal Component Analysis}
% ===========================================================================
Principal Component Analysis (PCA), pioneered independently by Karl Pearson \cite{Pearson_1901} and Harold Hotelling \cite{Hotelling_1933}, rotates the variables into a new orthogonal set of axes — the \emph{principal components} — ordered so that the first captures the most variance, the second the most of what remains, and so on. By transforming correlated variables into a smaller number of uncorrelated principal components, PCA effectively separates the signal from the noise in high-dimensional data. Following the revised schedule, this chapter also folds in \emph{Principal Coordinate Analysis} (PCoA), which we show is PCA seen through a distance matrix.

\section{Three routes to the same answer}
PCA can be computed as (i) the spectral decomposition of the correlation matrix
$\mathbf{R}=\mathbf{V}\mathbf{\Lambda}\mathbf{V}^{\!\top}$, (ii) the SVD of the scaled
data matrix, or (iii) \R's \code{princomp}/\code{prcomp}. All three return the same
eigenvalues (component variances), loadings (the $\mathbf{V}$ axes), and scores (the
rotated data $\mathbf{Z}=\mathbf{X}_s\mathbf{V}$). \code{prcomp} is preferred for
numerical stability because it uses SVD directly.

\begin{lstlisting}
Xs   <- scale(as.matrix(iris[, 1:4]))
R    <- cor(iris[, 1:4])
eig  <- eigen(R)                     # by hand
pc   <- prcomp(iris[, 1:4], scale. = TRUE)
rbind(eigenvalue   = eig$values,
      prop.var     = eig$values / sum(eig$values),
      cumulative   = cumsum(eig$values / sum(eig$values)))
\end{lstlisting}

\begin{lstlisting}[style=rout]
                  PC1      PC2      PC3      PC4
eigenvalue   2.918498 0.914030 0.146757 0.020715
prop.var     0.729624 0.228508 0.036689 0.005179
cumulative   0.729624 0.958132 0.994821 1.000000
\end{lstlisting}

Two components already carry $95.8\%$ of the variance. The equivalence of the routes is
not merely asserted — the code aligns signs and checks
$\max|\,\text{hand scores}-\text{\code{prcomp} scores}\,|\approx 10^{-15}$.

\subsection{Theory: PCA as a variance-maximization problem}
The first principal component is the unit vector $\mathbf{v}_1$ that maximises the
variance of the projected data $\mathbf{X}_s\mathbf{v}$. With $\mathbf{R}$ the
correlation matrix, this is
\[
  \mathbf{v}_1=\arg\max_{\|\mathbf v\|=1}\ \mathbf v^{\!\top}\mathbf R\,\mathbf v .
\]
Introducing a Lagrange multiplier $\lambda$ for the constraint $\mathbf v^{\!\top}\mathbf v=1$
and differentiating $\mathbf v^{\!\top}\mathbf R\mathbf v-\lambda(\mathbf v^{\!\top}\mathbf v-1)$
gives the eigen-equation
\[
  \mathbf R\,\mathbf v=\lambda\,\mathbf v,
\]
so the maximiser is the eigenvector of $\mathbf R$ with the largest eigenvalue, and that
eigenvalue \emph{is} the variance captured, $\mathbf v_1^{\!\top}\mathbf R\mathbf v_1=\lambda_1$.
Each subsequent component maximises the remaining variance subject to orthogonality with
the previous ones, yielding the full eigenbasis $\mathbf R=\mathbf V\mathbf\Lambda\mathbf V^{\!\top}$.
Because $\sum_j\lambda_j=\operatorname{tr}(\mathbf R)=p$ (here $p=4$), the proportion of
variance on component $j$ is simply $\lambda_j/p$ — the numbers in the table above.

\section{How many components? The scree plot}
The \emph{scree plot} graphs eigenvalue against component index; one keeps components
before the ``elbow.'' The \emph{Kaiser rule} keeps components with eigenvalue $>1$
(more variance than a single standardized variable). For iris both rules agree on two.

\begin{figure}[htbp]\centering
  \includegraphics[width=0.56\linewidth]{pca_scree.png}
  \caption{Iris PCA scree plot. Only the first two components exceed the Kaiser
  threshold (dashed line at eigenvalue $=1$).}
  \label{fig:pcascree}
\end{figure}

\section{The biplot}
A \emph{biplot} overlays object scores and variable loadings in one picture. Points are
objects in PC space; arrows are variables, their direction showing which components they
drive and their projection onto an axis showing their contribution.

\begin{figure}[htbp]\centering
  \includegraphics[width=0.66\linewidth]{pca_biplot.png}
  \caption{Iris PCA biplot. PC1 separates \emph{setosa} (left) from the others and is
  driven by the three positively-correlated size variables; sepal width points downward,
  loading on PC2.}
  \label{fig:pcabi}
\end{figure}

\section{PCoA: PCA through a distance matrix}
Principal Coordinate Analysis (PCoA), formalized by John Gower \cite{GOWER_1966}, takes a \emph{distance} matrix and finds
a set of Euclidean coordinates whose inter-point distances exactly reproduce the original dissimilarities — this is also known as classical (metric) multidimensional scaling. When the input distance is Euclidean on the same standardized data, PCoA returns the PCA scores exactly. This fundamental equivalence is why the standalone PCoA session is folded here: it is the exact same underlying eigen-problem, generalised to arbitrary non-Euclidean distances (like Gower or Bray--Curtis). This generalisation provides a powerful conceptual bridge, letting us ordinate mixed-type or count data that traditional PCA cannot handle.

\begin{lstlisting}
D    <- dist(scale(iris[, 1:4]))     # Euclidean distance
pcoa <- cmdscale(D, k = 2, eig = TRUE)
cor(pcoa$points[,1], prcomp(iris[,1:4], scale.=TRUE)$x[,1])   # ~ 1.000
\end{lstlisting}

\begin{figure}[htbp]\centering
  \includegraphics[width=0.56\linewidth]{09_pcoa_vs_pca.png}
  \caption{PCoA on a Euclidean distance reproduces the PCA ordination point for point.
  For non-Euclidean distances PCoA generalises PCA to data types PCA cannot handle.}
  \label{fig:pcoa}
\end{figure}

\begin{pitfall}
PCA on the \emph{covariance} matrix lets high-variance variables dominate; on the
\emph{correlation} matrix (scaled data) it treats variables even-handedly. State which
you used — they can give qualitatively different ordinations.
\end{pitfall}

% ===========================================================================
\chapter{Correspondence Analysis}
% ===========================================================================
Correspondence Analysis (CA) is PCA adapted to \emph{counts} — contingency tables of
non-negative integers, such as species abundances across sites. Instead of Euclidean
distance it uses the $\chi^2$ distance, and it ordinates rows (sites) and columns
(species) \emph{simultaneously}, because it is the SVD of a single $\chi^2$-standardized
matrix (recall the Q/R eigenvalue identity from Chapter~1).

\section{The mechanics}
Let $\mathbf{P}=\mathbf{N}/n_{\!++}$ be the table of relative frequencies with row and
column masses $\mathbf{r},\mathbf{c}$. CA takes the SVD of
$\mathbf{S}=\mathbf{D}_r^{-1/2}(\mathbf{P}-\mathbf{r}\mathbf{c}^{\!\top})\mathbf{D}_c^{-1/2}$.
The squared singular values are the \emph{inertias} (variance explained per axis).

\begin{lstlisting}
library(vegan)
fish <- get_data("doubs")$fish      # 30 sites x 27 fish species (counts)
ca   <- cca(fish)                    # correspondence analysis in vegan
round(ca$CA$eig / sum(ca$CA$eig), 3)[1:4]
\end{lstlisting}

\begin{lstlisting}[style=rout]
  CA1   CA2   CA3   CA4
0.317 0.241 0.159 0.085
\end{lstlisting}

\begin{figure}[htbp]\centering
  \includegraphics[width=0.62\linewidth]{ca_biplot.png}
  \caption{CA of a river fish gradient. Sites (points, coloured by upstream--downstream
  position) and species (triangles) share the same space; the first axis recovers the
  longitudinal gradient of the river.}
  \label{fig:ca}
\end{figure}

\begin{pitfall}
Long ecological gradients induce the \emph{arch (horseshoe) effect} in CA — a curved
smile with no real second gradient. Detrended CA (\code{decorana}) removes it, but treat
detrending as a diagnostic, not a default.
\end{pitfall}

% ===========================================================================
\chapter{Multidimensional Scaling}
% ===========================================================================
Where PCA and CA are eigen-methods, \emph{non-metric multidimensional scaling} (nMDS)
is an \emph{optimisation}. It seeks a low-dimensional configuration whose inter-point
distances preserve the \emph{rank order} of the original dissimilarities, minimising a
badness-of-fit criterion called \emph{stress}. It makes no distributional assumptions and
accepts any dissimilarity, which makes it the most flexible ordination in the toolbox.

\begin{lstlisting}
library(vegan)
md <- metaMDS(data, distance = "bray", k = 2)   # random-start optimisation
md$stress                                        # < 0.1 good, < 0.2 usable
stressplot(md)                                   # Shepard diagram
\end{lstlisting}

\begin{figure}[htbp]\centering
  \includegraphics[width=\linewidth]{mds_shepard.png}
  \caption{Left: an nMDS configuration. Right: the Shepard diagram plots configuration
  distance against original distance; a tight, monotone increasing cloud (near the red
  line) means the low-dimensional map faithfully preserves the dissimilarity ranks.}
  \label{fig:mds}
\end{figure}

\begin{keybox}
Stress is an optimisation score, not a probability. A low stress means the picture is a
faithful map of the distances; it does \emph{not} certify that the distance you chose was
the right question. Garbage distance in, low-stress garbage out.
\end{keybox}

% ===========================================================================
\chapter{Non-negative Matrix Factorization}
% ===========================================================================
Non-negative Matrix Factorization (NMF) approximates a non-negative data matrix as
$\mathbf{A}\approx\mathbf{W}\mathbf{H}$ with \emph{both} factors non-negative. Because
nothing can cancel, the columns of $\mathbf{W}$ act as additive ``parts'' (metagenes,
motifs, topics) and the columns of $\mathbf{H}$ are the amounts of each part in each
sample. The non-negativity buys interpretability that PCA's signed loadings lack.

\begin{lstlisting}
library(NMF)
expr <- get_data("leukemia")         # 300 genes x 60 patients (3 subtypes)
res  <- nmf(expr - min(expr), rank = 3, seed = "nndsvd")
basismap(res); coefmap(res)          # W (gene signatures) and H (sample loadings)
\end{lstlisting}

\begin{figure}[htbp]\centering
  \includegraphics[width=\linewidth]{nmf_heatmaps.png}
  \caption{Rank-3 NMF of a leukemia expression matrix. Left: the basis $\mathbf{W}$
  (three gene signatures). Right: the coefficients $\mathbf{H}$; the three factors line
  up with the three clinical subtypes (white separators), recovering biology unsupervised.}
  \label{fig:nmf}
\end{figure}

% ===========================================================================
\chapter{Exploratory Factor Analysis}
% ===========================================================================
Factor analysis (FA) looks superficially like PCA but rests on a different model. PCA
re-expresses \emph{all} the variance with components that are exact linear combinations
of the variables. FA posits a small number of unobserved \emph{common factors} that
generate the correlations among variables, and it explicitly separates shared variance
from variable-specific noise.

\section{The model}
With $p$ standardized variables and $k<p$ latent factors $\mathbf f$,
\[
  \mathbf x=\mathbf\Lambda\,\mathbf f+\boldsymbol\varepsilon,\qquad
  \operatorname{Cov}(\mathbf x)=\mathbf\Lambda\mathbf\Lambda^{\!\top}+\mathbf\Psi,
\]
where $\mathbf\Lambda$ ($p\times k$) holds the \emph{loadings}, $\mathbf f$ are the
common factors (uncorrelated, unit variance), and $\mathbf\Psi$ is a diagonal matrix of
\emph{uniquenesses} — the part of each variable not explained by the common factors. The
key contrast with PCA is that term $\mathbf\Psi$: FA does not try to fit variable-specific
variance, only the \emph{common} structure. Estimation is by maximum likelihood, and the
loadings are then \emph{rotated} (e.g.\ varimax) to a more interpretable, near-simple
structure.

\begin{lstlisting}
X  <- scale(get_data("taxon")[, -1])
fa <- factanal(X, factors = 2, rotation = "varimax")
round(eigen(cor(X))$values, 3)          # how many factors?
print(fa$loadings, digits = 2, cutoff = 0.2)
\end{lstlisting}

\begin{lstlisting}[style=rout]
correlation-matrix eigenvalues:
[1] 4.311 1.512 0.377 0.295 0.232 0.151 0.123     # two exceed 1 (Kaiser)

Loadings (varimax):
          Factor1 Factor2
Petals     -0.47    0.69
Internode   .      -0.90
Sepal       0.77   -0.35
Bract      -0.69    0.39
Petiole    -0.84    .
Leaf        .      -0.90
Fruit       0.88   -0.36
\end{lstlisting}

Two eigenvalues exceed one, so two factors are retained. After varimax rotation each
variable loads strongly on one factor and weakly on the other — the ``simple structure''
that makes factors interpretable. (The correlation-matrix eigenvalues above are exact;
the loading \emph{signs and magnitudes} depend on the estimator and rotation, so read the
\emph{pattern} of high/low loadings rather than individual decimals.)

\begin{figure}[htbp]\centering
  \includegraphics[width=0.5\linewidth]{factor_scree.png}
  \caption{Scree of the correlation-matrix eigenvalues; the Kaiser line (eigenvalue $=1$)
  retains two factors, consistent with the four well-separated groups in the data.}
  \label{fig:fa}
\end{figure}

\begin{pitfall}
PCA and FA are not interchangeable. Use PCA to \emph{reduce} dimensionality (summarise
total variance); use FA to \emph{explain} correlations by latent constructs. Reporting
PCA loadings as ``factors'' — a common slip — conflates the two models.
\end{pitfall}

% ###########################################################################
\part{Classification and Association}
% ###########################################################################

% ===========================================================================
\chapter{Clustering}
% ===========================================================================
Clustering partitions objects into groups without using labels. We cover the two
classroom workhorses: \emph{k-means} (partition around $k$ centroids) and
\emph{hierarchical} clustering (a nested tree, or dendrogram).

\section{k-means and the elbow}
k-means minimises the total within-cluster sum of squares (WSS). Because WSS always
decreases with $k$, we look for the \emph{elbow} where added clusters stop paying off.

\begin{lstlisting}
w <- wss(as.matrix(iris[,1:4]))      # WSS for k = 1..10 (helper in 05_Clustering.R)
plot(w, type = "b"); abline(v = 3, lty = 2)
km3 <- kmeans(iris[,1:4], 3, nstart = 25)
table(cluster = km3$cluster, species = iris$Species)
\end{lstlisting}

\begin{lstlisting}[style=rout]
        species
cluster  setosa versicolor virginica
      1      50          0         0
      2       0         48        14
      3       0          2        36
\end{lstlisting}

The elbow sits at $k=3$ (Figure~\ref{fig:elbow}). The cross-tabulation shows the
unsupervised clusters recover the species with $\approx 89\%$ agreement: \emph{setosa}
is perfect, while \emph{versicolor}/\emph{virginica} overlap — exactly the structure the
Q-mode heatmap warned about in Chapter~1.

\begin{figure}[htbp]\centering
  \includegraphics[width=0.52\linewidth]{kmeans_elbow.png}\hfill
  \includegraphics[width=0.44\linewidth]{kmeans_crosstab.png}
  \caption{Left: k-means elbow — WSS drops sharply to $k=3$ then flattens. Right: the
  $k=3$ clusters versus true species; only 16 of 150 objects are ``misplaced,'' all at
  the versicolor/virginica boundary.}
  \label{fig:elbow}
\end{figure}

\section{Hierarchical clustering and linkage}
Agglomerative clustering repeatedly merges the two closest groups. The \emph{linkage}
rule defines ``closest'': single (nearest point, chains), complete (farthest point,
compact), average (UPGMA), and Ward (minimum variance, the ecologist's default).

\begin{figure}[htbp]\centering
  \includegraphics[width=0.7\linewidth]{hclust_dendrogram.png}
  \caption{Ward dendrogram of world regions clustered on standardized biodiversity
  variables. Cutting the tree at a chosen height yields discrete groups.}
  \label{fig:dendro}
\end{figure}

\begin{pitfall}
k-means depends on the random initial centroids; always use \code{nstart} $>1$ (we use
25) so the reported solution is the best of many restarts, not a lucky or unlucky one.
Different linkages can produce very different trees from the \emph{same} distances — the
choice is substantive, not cosmetic.
\end{pitfall}

% ===========================================================================
\chapter{Discriminant Analysis}
% ===========================================================================
Discriminant analysis is \emph{supervised} ordination: given labelled groups, it finds
the axes that best \emph{separate} them. Linear Discriminant Analysis (LDA) maximises the
ratio of between-group to within-group variance, assuming equal group covariances;
Quadratic DA (QDA) relaxes that assumption.

\subsection{Theory: Fisher's criterion}
Let $\mathbf B$ be the between-group and $\mathbf W$ the within-group scatter matrix.
Fisher sought the direction $\mathbf a$ maximising the separation ratio
\[
  J(\mathbf a)=\frac{\mathbf a^{\!\top}\mathbf B\,\mathbf a}{\mathbf a^{\!\top}\mathbf W\,\mathbf a}.
\]
Setting $\partial J/\partial\mathbf a=\mathbf 0$ yields the \emph{generalized}
eigenproblem $\mathbf B\mathbf a=\ell\,\mathbf W\mathbf a$, i.e.\ the ordinary eigenproblem
$\mathbf W^{-1}\mathbf B\,\mathbf a=\ell\,\mathbf a$. The leading eigenvectors are the
discriminant axes and their eigenvalues (the ``proportion of trace'') say how much group
separation each axis carries. With $g$ groups there are at most $g-1$ such axes — here,
three for four taxa.

\begin{lstlisting}
library(MASS)
model <- lda(Taxon ~ ., data = train)      # taxon: 4 morphometric groups
pred  <- predict(model, test)$class
mean(pred == test$Taxon)                    # hold-out accuracy
\end{lstlisting}

\begin{lstlisting}[style=rout]
[1] 1.000       # perfect separation on held-out data (well-separated groups)
\end{lstlisting}

Figure~\ref{fig:lda} makes the point of supervision vivid: PCA, which never sees the
labels, smears the four groups along directions of raw variance; LDA, which uses them,
opens clean gaps. That gain is real when groups differ, and illusory when they do not —
hence the assumption checks below.

\begin{figure}[htbp]\centering
  \includegraphics[width=\linewidth]{lda_vs_pca.png}
  \caption{The same four-group data under PCA (left, unsupervised) and LDA (right,
  supervised). LDA's axes are built to separate the known classes.}
  \label{fig:lda}
\end{figure}

\begin{pitfall}
LDA assumes multivariate normality and \emph{equal covariance matrices} across groups.
Test the second assumption (e.g.\ \code{vegan::betadisper} with a permutation test)
before trusting an LDA; when it fails, prefer QDA or a tree-based classifier from the
next chapter.
\end{pitfall}

% ===========================================================================
\chapter{Multivariate Generalized Linear Models}
% ===========================================================================
Ordinary regression assumes a continuous, unbounded response with Gaussian, constant-
variance errors. Ecological responses rarely oblige: biomass is positive, counts are
integers, presence is binary. \emph{Generalized linear models} (GLMs) keep the linear
predictor $\eta=\mathbf x^{\!\top}\boldsymbol\beta$ but connect it to the mean through a
\emph{link} function and allow a non-Gaussian error \emph{family}.

\section{Link and family}
A GLM has three parts: a linear predictor $\eta$; a link $g$ with $g(\mu)=\eta$; and an
exponential-family distribution for the response. Choose the family by the response type
— Gamma (log link) for positive continuous data, Poisson (log link) for counts, binomial
(logit link) for 0/1 — and fit by iteratively reweighted least squares.

\begin{lstlisting}
m <- get_data("crawley")            # mixed-type ecological table
shapiro.test(m$Biomass)$p.value     # small p => Gaussian OLS is wrong
g_gamma <- glm(Biomass ~ Tmp + pH, family = Gamma(link = "log"), data = m)
g_pois  <- glm(Species ~ Tmp + pH, family = poisson,             data = m)
AIC(g_pois, glm(Species ~ Tmp * pH, family = poisson, data = m))  # additive vs interaction
\end{lstlisting}

\begin{lstlisting}[style=rout]
Shapiro-Wilk on Biomass: p-value < 1e-6      # reject normality -> use a GLM

Gamma GLM (Biomass ~ Tmp + pH), log link:
(Intercept)         Tmp          pH
     1.6070      0.0673     -0.0415       # on the log scale

Poisson GLM (Species ~ Tmp + pH) coefficients:
(Intercept)         Tmp          pH
     0.5028      0.0543      0.1569       # on the log-count scale

AIC additive vs interaction (Poisson):
              df      AIC
g_pois         3   404.7
g_int          4   406.7              # interaction not worth the extra parameter
\end{lstlisting}

The Shapiro test rejects normality of the (right-skewed) biomass, justifying the GLM; the
coefficients live on the log scale, so $e^{\beta}$ is a multiplicative effect on the mean;
and AIC prefers the simpler additive model (adding the interaction raises AIC from
$404.7$ to $406.7$).

\begin{figure}[htbp]\centering
  \includegraphics[width=\linewidth]{19_glm.png}
  \caption{GLM fits: a Gamma model for positive-continuous biomass (left) and a Poisson
  model for species counts (right). The link keeps predictions in the valid range that a
  straight OLS line would violate.}
  \label{fig:glm}
\end{figure}

\begin{pitfall}
Poisson regression assumes mean $=$ variance. Real counts are often \emph{overdispersed}
(variance $>$ mean); then use a quasi-Poisson or negative-binomial family, or the standard
errors will be too small and everything looks significant.
\end{pitfall}

% ===========================================================================
\chapter{Canonical Correlation and Redundancy Analysis}
% ===========================================================================
The final classical methods generalise correlation and regression from vectors to
\emph{two tables}. Canonical Correlation Analysis (CCorA) is the two-table analogue of
correlation — symmetric, no direction. Redundancy Analysis (RDA) is the two-table
analogue of regression — directional, one table explained by the other.

\section{Canonical correlation}
Given tables $\mathbf X$ ($n\times p$) and $\mathbf Y$ ($n\times q$), CCorA finds pairs of
linear combinations $\mathbf X\mathbf a$ and $\mathbf Y\mathbf b$ with maximal
correlation, then a second orthogonal pair, and so on. The canonical correlations are the
square roots of the eigenvalues of
$\mathbf R_{xx}^{-1}\mathbf R_{xy}\mathbf R_{yy}^{-1}\mathbf R_{yx}$, and Pillai's trace
(with a permutation test) assesses whether the two tables are associated at all.

\begin{lstlisting}
doubs <- get_data("doubs")
Y <- decostand(doubs$fish, "hellinger")    # response table (fish)
X <- scale(doubs$env)                       # predictor table (environment)
CCorA(scale(doubs$fish), X)$p.perm          # are the tables associated?
rda.mod <- rda(Y ~ ., data = as.data.frame(X))   # directional: fish ~ environment
RsquareAdj(rda.mod)$adj.r.squared
\end{lstlisting}

\section{Redundancy analysis}
RDA regresses \emph{every} response variable on the predictors, collects the fitted
values into a matrix, and runs a PCA on that fitted matrix. The variance thus splits into
a \emph{constrained} part (explained by $\mathbf X$) and an \emph{unconstrained} part
(residual), and an \emph{adjusted} $R^2$ corrects the optimistic raw fraction for the
number of predictors. A permutation test (\code{anova.cca}) provides overall
significance, and \code{ordiR2step} does forward selection of predictors.

\begin{figure}[htbp]\centering
  \includegraphics[width=0.66\linewidth]{20_rda.png}
  \caption{RDA triplot: sites (points), species (text) and environmental predictors
  (arrows). Arrow length and direction show which gradients drive community composition;
  this is multivariate multiple regression made visual.}
  \label{fig:rda}
\end{figure}

\begin{keybox}
Canonical correlation asks ``are these two tables related?'' (symmetric); redundancy
analysis asks ``how much of table $\mathbf Y$ does table $\mathbf X$ explain?''
(directional). RDA is the multivariate regression that closes the classical half of the
course — and it partitions variance exactly as an ANOVA does, but for a whole community
at once.
\end{keybox}

% ###########################################################################
\part{Machine Learning}
% ###########################################################################

% ===========================================================================
\chapter{Supervised Machine Learning}
% ===========================================================================
This chapter integrates the \pkg{CART} and \pkg{caret} units of the machine-learning
course. The classical classifiers of Chapter~8 assume a functional form; tree-based
learners instead \emph{partition} the feature space with axis-aligned splits, capturing
interactions and non-linearities with no distributional assumptions.

\section{From one tree to a forest}
A single decision tree (CART) is interpretable but high-variance. Three ensemble ideas
tame that variance: \emph{bagging} averages trees on bootstrap samples; the
\emph{random forest} additionally de-correlates them by sampling variables at each split;
\emph{boosting} grows trees sequentially, each correcting its predecessor.

\begin{notebox}
The code here is lifted, lightly adapted, from the \code{UnitCART} and \code{UnitCARET}
scripts of the ML course (both graded S+). The pedagogical move is to present trees and
forests as the supervised cousins of the clustering and ordination methods students have
already met — same data matrix, new objective.
\end{notebox}

\begin{lstlisting}
library(rpart); library(randomForest); library(caret)
set.seed(1998)
idx   <- createDataPartition(iris$Species, p = 0.7, list = FALSE)
train <- iris[idx, ]; test <- iris[-idx, ]

fit_cart <- rpart(Species ~ ., data = train, method = "class")
fit_rf   <- randomForest(Species ~ ., data = train, ntree = 500, importance = TRUE)

## one interface, honest 5-fold CV, automatic tuning
ctrl <- trainControl(method = "cv", number = 5)
m_rf <- train(Species ~ ., data = train, method = "rf", trControl = ctrl)
\end{lstlisting}

\begin{lstlisting}[style=rout]
5-fold cross-validated accuracy (mean):
  CART (rpart)      0.960
  Random Forest     0.960
  Gradient Boosting 0.960
\end{lstlisting}

On iris all three reach $\approx0.96$; the value of the forest shows on harder, noisier
problems, where a single tree overfits. Figure~\ref{fig:mlbound} shows \emph{why} the
methods differ: the single tree carves the plane into a few rectangles, the forest
smooths those steps by averaging, and boosting sharpens the class boundaries.

\begin{figure}[htbp]\centering
  \includegraphics[width=\linewidth]{ml_decision_boundaries.png}
  \caption{Decision boundaries on the two petal variables. The single CART tree (left)
  uses a few axis-aligned cuts; the random forest (middle) and gradient boosting (right)
  produce smoother, more robust partitions.}
  \label{fig:mlbound}
\end{figure}

\begin{figure}[htbp]\centering
  \includegraphics[width=0.52\linewidth]{ml_rf_importance.png}\hfill
  \includegraphics[width=0.44\linewidth]{ml_model_comparison.png}
  \caption{Left: random-forest variable importance — the petal measurements dominate,
  echoing the PCA loadings of Chapter~3. Right: cross-validated accuracy of the three
  tuned models.}
  \label{fig:mlimp}
\end{figure}

\begin{keybox}
The workflow, not the algorithm, is the lesson: split the data, resample honestly
(cross-validation), tune on the training folds only, and report on a held-out test set.
\pkg{caret} (and its successor \pkg{tidymodels}) makes this one consistent interface
across dozens of models.
\end{keybox}

% ===========================================================================
\chapter{Unsupervised Machine Learning: t-SNE and UMAP}
% ===========================================================================
PCA is a \emph{linear} projection: it preserves global variance but can leave curved,
clustered structure tangled. Modern \emph{manifold learners} — t-SNE and UMAP — are
non-linear and optimise for preserving \emph{local} neighbourhoods, so they pull apart
clusters that PCA overlaps. They are the unsupervised counterpart to the ordinations of
Part II, and the standard first look at high-dimensional omics data.

\begin{lstlisting}
library(Rtsne); library(uwot)
X  <- t(get_data("leukemia"))         # 60 patients x 300 genes
pc <- prcomp(X, scale. = TRUE)$x[, 1:2]
ts <- Rtsne(X, perplexity = 15)$Y     # t-SNE
um <- umap(X, n_neighbors = 15)       # UMAP
\end{lstlisting}

\begin{figure}[htbp]\centering
  \includegraphics[width=\linewidth]{unsup_embeddings.png}
  \caption{Three views of the same leukemia expression matrix. PCA (left) already hints
  at the three subtypes; t-SNE (middle) and UMAP (right) sharpen them into clearly
  isolated islands by honouring local neighbourhoods.}
  \label{fig:embed}
\end{figure}

\begin{pitfall}
t-SNE and UMAP distort \emph{global} geometry: distances \emph{between} well-separated
clusters, and cluster \emph{sizes}, are not meaningful, and the layout depends on
perplexity / \code{n\_neighbors}. Use them to \emph{see} groups, not to measure how far
apart groups are — for that, return to PCA or PCoA.
\end{pitfall}

% ===========================================================================
\chapter{Neural Networks and Deep Learning}
% ===========================================================================
A neural network stacks linear maps and non-linear activations so that, with enough
hidden units, it can approximate essentially any function. Training adjusts the weights
by gradient descent to minimise a loss. This chapter integrates the \code{UnitNN} deep-
learning unit of the ML course.

\section{A network you can run anywhere}
The self-contained track uses base-\R's \pkg{nnet}: one hidden layer, softmax output,
weight decay for regularisation. It needs no Python, keras or GPU, so it always runs.

\begin{lstlisting}
library(nnet)
Xs  <- scale(as.matrix(iris[,1:4]))
Y   <- class.ind(iris$Species)              # one-hot targets
net <- nnet(Xs[idx,], Y[idx,], size = 8, softmax = TRUE, decay = 5e-3, maxit = 400)
mean(predict(net, Xs[-idx,], type = "class") == iris$Species[-idx])
\end{lstlisting}

\begin{lstlisting}[style=rout]
[1] 0.98    # representative held-out accuracy (0.96-1.00 across random splits/
            # weight inits); architecture 4 -> 8 -> 3, 67 weights
\end{lstlisting}

\begin{figure}[htbp]\centering
  \includegraphics[width=0.52\linewidth]{nn_training_curve.png}\hfill
  \includegraphics[width=0.44\linewidth]{21_nn_boundary.png}
  \caption{Left: training and validation error fall together as the network learns
  (the validation curve is the one to watch for overfitting). Right: the network's
  non-linear decision surface on the two petal variables.}
  \label{fig:nn}
\end{figure}

\begin{notebox}
The chapter's second track is a \pkg{keras}/\pkg{tensorflow} template — a dense classifier
and a small convolutional network — copied from \code{UnitNN} and guarded by
\code{requireNamespace()} so the script never breaks when keras is absent. Students with a
GPU get the full deep-learning workflow (image track included); everyone else still runs
the \pkg{nnet} version and sees the same concepts.
\end{notebox}

\begin{pitfall}
Neural networks are the most eager over-fitters in this book. They demand standardized
inputs, explicit regularisation (weight decay, dropout), and an honest validation split.
On the small tabular data typical of ecology, a random forest usually matches or beats a
neural net with a fraction of the tuning.
\end{pitfall}

% ===========================================================================
\chapter{Large Language Models and Embeddings}
% ===========================================================================
The course ends where machine learning meets the theme of Chapter~1. A large language
model (LLM) turns each word, sentence or document into a high-dimensional vector — an
\emph{embedding} — arranged so that geometric proximity encodes \emph{meaning}. An
embedding \emph{is} a row of a multivariate data matrix, so every tool in this book
applies to it unchanged: distances, PCA/UMAP, clustering, classification.

\section{Meaning as geometry}
For unit-normalised embeddings the natural similarity is the cosine (the dot product).
Remarkably, \emph{relations} become vector arithmetic: the direction from \code{man} to
\code{woman} parallels \code{king} to \code{queen}, so
$\text{king}-\text{man}+\text{woman}\approx\text{queen}$.

\begin{lstlisting}
E <- <embedding matrix, one unit-row per word>   # from an LLM, or built structurally
q <- E["king",] - E["man",] + E["woman",]
sort(E %*% (q / sqrt(sum(q^2))), decreasing = TRUE)[1:3]
\end{lstlisting}

\begin{lstlisting}[style=rout]
 queen  woman  italy      # nearest neighbours to (king - man + woman)
 0.83   0.61   0.19
\end{lstlisting}

\begin{figure}[htbp]\centering
  \includegraphics[width=0.62\linewidth]{llm_embeddings.png}
  \caption{A PCA projection of word embeddings. Semantically related words cluster —
  royalty/gender, geography, animals, biology — because the embedding places meaning in
  geometry. The same PCA, distances and clustering from earlier chapters now operate on
  language.}
  \label{fig:llm}
\end{figure}

\begin{keybox}
Nothing in this chapter is new machinery — it is the \emph{data matrix} of Chapter~1,
populated by a language model instead of a caliper. Transformers (Vaswani et al., 2017)
produce the vectors; the multivariate toolbox of this book interprets them. That is the
throughline of the course: master the geometry of the data matrix, and every method,
classical or modern, becomes a variation on one theme.
\end{keybox}

% ###########################################################################
\part{Applications in Spatial Ecology}
% ###########################################################################

% ===========================================================================
\chapter{Ecological Niche Modeling}
% ===========================================================================
An ecological niche model (ENM) estimates the environmental conditions where a species
can occur. Statistically it is a \emph{binomial regression} — the GLM of Chapter~11 —
whose response is presence/absence (or presence/background) and whose predictors are
climate variables. The fitted surface, projected onto geography, becomes a suitability
map. The full classroom pipeline uses spatial rasters (\pkg{terra}, \pkg{dismo},
WorldClim); here we distil the statistical core so it runs without GIS layers.

\begin{lstlisting}
occ  <- get_data("hanta")               # Sp (0/1), bio_1, bio_12
enm2 <- glm(Sp ~ bio_1 + bio_12, data = occ, family = binomial)
enmq <- glm(Sp ~ poly(bio_1, 2) + bio_12, data = occ, family = binomial)
AIC(enm2, enmq)                         # does a hump-shaped temperature response help?
\end{lstlisting}

\begin{lstlisting}[style=rout]
Best-model coefficients (bio_1 + bio_12):
(Intercept)       bio_1      bio_12
   -2.15504     0.13201    -0.00152

Training AUC (bio_1 + bio_12): 0.843     # good discrimination (0.5 = random, 1 = perfect)
\end{lstlisting}

The positive \code{bio\_1} coefficient means warmer sites are more suitable; the small
negative \code{bio\_12} term slightly penalises the wettest sites. An AUC of $0.843$ says
the model ranks a random presence above a random absence about $84\%$ of the time.
Figure~\ref{fig:enm} shows the two standard ENM outputs: a \emph{response curve} (how
suitability varies along one predictor) and a \emph{suitability surface} in environmental
space, with observed detections overlaid.

\begin{figure}[htbp]\centering
  \includegraphics[width=\linewidth]{enm.png}
  \caption{Ecological niche model. Left: response curve of presence probability against
  temperature. Right: the fitted suitability surface over the two bioclimatic axes;
  detections (points) fall in the high-suitability region.}
  \label{fig:enm}
\end{figure}

\begin{pitfall}
Presence-only data has no true absences; the ``0''s are usually \emph{background} points.
Treating background as genuine absence biases prevalence and inflates apparent
performance. And AUC computed on the training data (as here, for illustration) is
optimistic — always report AUC on spatially independent held-out data.
\end{pitfall}

% ===========================================================================
\chapter{Presence--Absence Matrices}
% ===========================================================================
A presence--absence matrix (PAM) records which species occur at which sites as a binary
matrix — the discrete cousin of the community tables from Correspondence Analysis. Its
two marginal sums are the fundamental quantities of biodiversity: row sums are
\emph{richness} (species per site, $\alpha$) and column sums are \emph{range size}
(sites per species, $\omega$).

\begin{lstlisting}
pamdf <- get_data("pam")
pam   <- as.matrix(pamdf[, -1])         # sites x species (0/1)
alpha <- rowSums(pam)                    # richness per site
omega <- colSums(pam)                    # range size per species
c(fill = mean(pam), mean_alpha = mean(alpha), mean_omega = mean(omega))
\end{lstlisting}

\begin{lstlisting}[style=rout]
       fill  mean_alpha  mean_omega
      0.327       39.27       26.18       # 80 sites x 120 species, 33% filled
\end{lstlisting}

\begin{figure}[htbp]\centering
  \includegraphics[width=\linewidth]{pam.png}
  \caption{A presence--absence matrix (left, sites ordered by latitude), the richness
  gradient it implies ($\alpha$ per site, middle), and the frequency distribution of
  range sizes ($\omega$ per species, right).}
  \label{fig:pam}
\end{figure}

\begin{keybox}
The matrix product $\mathbf{P}\mathbf{P}^{\!\top}$ counts \emph{shared species} between
every pair of sites, and $\mathbf{P}^{\!\top}\mathbf{P}$ counts \emph{co-occurrences}
between every pair of species. From these two products fall out $\beta$-diversity, dispersion
fields, and range--diversity plots — the whole macroecology of the PAM is linear algebra
on a binary matrix, closing the loop back to Chapter~1.
\end{keybox}

% ###########################################################################
\appendix
% ###########################################################################
\chapter{Reproducing this book}
Everything here is generated from the \code{MS\_LJMR/} package.

\textbf{Layout.}
\code{R/} holds the revised, portable course scripts (no \code{setwd()}, no \code{x11()};
data obtained through \code{get\_data()}). \code{verify/} holds the Python reproductions
that cross-check every method and render the figures. \code{figs/} holds those figures.
\code{book/} holds this single \code{.tex} file. \code{data/} is where the real course
data sets go — when present they are used automatically.

\textbf{To rebuild the figures:}
\begin{lstlisting}
cd MS_LJMR/verify
python3 make_figures.py      # core chapters + verification log
python3 extra_figures.py     # dissimilarity, correlation, MDS figures
\end{lstlisting}

\textbf{To rebuild this book:}
\begin{lstlisting}
cd MS_LJMR/book
pdflatex MS_LJMR_Multivariate.tex   # run 2-3 times for the table of contents
\end{lstlisting}

\textbf{Verification.} Because no \R{} interpreter was available while assembling this
edition, each method's key quantities were recomputed independently in Python
(NumPy / SciPy / scikit-learn / statsmodels) and logged to \code{verify/verify\_log.txt};
the \R{} scripts were written to match. Representative reference values: iris PCA variance
$73.0/22.9/3.7/0.5\%$; $k$-means recovers species with $89\%$ agreement; LDA hold-out
accuracy $1.00$ on the four-group morphometrics; tree/forest/boosting $\approx0.96$
5-fold accuracy; the \code{king$-$man$+$woman$\to$queen} analogy resolves correctly.

\chapter{Notes on the revision}
This edition follows the revised Fall-2026 schedule: Principal Coordinate Analysis is
folded into the PCA chapter (Chapter~3); the standalone visualization and \pkg{ggplot2}
sessions are retired (their techniques are used throughout); and four machine-learning
chapters (9--12) are added, adapted from the author's graduate ML course. Bugs in the
original scripts were fixed in passing — an undefined \code{d\_euc2} in the dissimilarity
script, a \code{Qroo3}/\code{QRoo3} name clash in the MDS-biplot script, and pervasive
hard-coded absolute paths — and every script now runs unmodified on any platform.

\vspace{1em}
{\small\itshape Selected references.} Gower (1966) \emph{Biometrika}; Legendre \&
Legendre (2012) \emph{Numerical Ecology}; James, Witten, Hastie \& Tibshirani (2021)
\emph{ISLR}; Cutler et al.\ (2007) \emph{Ecology}; van der Maaten \& Hinton (2008)
\emph{JMLR}; McInnes et al.\ (2018) UMAP; Vaswani et al.\ (2017) \emph{NeurIPS};
Borowiec et al.\ (2022) \emph{Methods Ecol.\ Evol.}

\bibliographystyle{plain}
\bibliography{references}

\end{document}
